\section{Theory}
This section will introduce all of the relevant theory for this thesis. 

\subsection{Compound nucleus mechanism}
Bohr suggested this in 1936. 

\subsection{R-matrix theory}
There are two main issues with the compound nucleus description: firstly, it strictly defines a specific physical picture of how a reaction proceeds \note{expand somewhere, from first page of Lane}, and secondly, it is based on time-dependent perturbation theory. It is especially hard to justify the second issue, since nuclear forces can certainly not be treated as mere perturbations. Thus R-matrix theory was developed to solve these issues. \\

Note that there is nothing wrong with the physical picture provided by the compound nucleus description. In fact, it is rather useful, and will be used here when describing R-matrix theory as well. The problem is that in the compound nucleus description, it is the \textit{only} picture available; R-matrix theory can theoretically describe all types of reaction mechanisms. \\

The goal is to describe the intermediate compound nucleus in terms of a set of states, or \textit{channels}, as they are called within the theory. To do this, an interaction radius $a_c$ is defined. The idea is that this radius divides the space between the two particles into two regions: outside the radius, the only interaction is the Coulomb force, while inside the radius the nuclear forces dominate. Since this radius obviously depends on the strength of the Coulomb force, it will in general depend on the interacting nuclei. A popular choice for its value is to use the nuclear radii \note{reference Lane p. 266 eq 1.1}: \begin{align*}
	a_c = r_0 (A^{1/3}_1 + A^{1/3}_2) && r_0 \approx 1.4fm
\end{align*}